\begin{abstract}
Living organisms maintain complex shapes against gravity, but the physical limits of this maintenance are unknown. Here we derive a dimensionless scaling parameter --- the Bio-Gravitational Number $\mathcal{B}_g = EI/(MgL^2)$ --- that predicts when passive structure becomes insufficient and active metabolic maintenance is required. Cross-species analysis confirms that humans uniquely occupy an ``Allometric Trap'' ($\mathcal{B}_g \approx 0.01$) where shape maintenance depends entirely on metabolic energy. Using an Information--Elasticity Coupling framework that connects HOX-patterned developmental fields to Cosserat rod mechanics, we show that the spinal S-curve emerges as the energetic ground state, reproducing clinical lordosis (${\sim}42^\circ$) and kyphosis (${\sim}35^\circ$). AlphaFold structural analysis of 23 mechanotransduction proteins reveals a 72\% demand--supply anisotropy gap ($p = 0.011$): the proteins required for gravity sensing are thermodynamically expensive, while their metabolic regulators are intrinsically disordered and fragile. This molecular asymmetry creates an Energy Deficit Window during adolescent growth ($L > 0.35$\,m) that our model predicts produces scoliosis-like instabilities. These results reclassify Adolescent Idiopathic Scoliosis as a thermodynamic scaling failure and identify molecular targets for metabolic intervention.
\end{abstract}
